Several groups have bumped into the validation gap without mapping it.
Robomimic~\citep{mandlekar2021} reported the 50-to-100-percent discrepancy in a discussion
section. ACT~\citep{zhao2023} used validation loss as a working default. SIMPLER
\citep{li2024} measured weak correlation between validation error and real success, then
turned to simulation-based evaluation rather than offline metrics. NeME~\citep{tiezzi2025}
trained one alternative scorer for humanoid interaction. A line of work on rollout-based
evaluation~\citep{vincent2024,snyder2026,vincent2026golden,anwar2025} accepts rollout cost
and makes it efficient, sidestepping the offline question entirely.

Two recent efforts come closest and must be distinguished. RoboEval~\citep{roboeval2025}
augments binary success with richer behavioural metrics (coordination, trajectory smoothness,
spatial precision) and shows they correlate with success across many tasks --- but those
metrics are measured \emph{during closed-loop execution}. We ask the opposite question:
whether any signal computable \emph{without} a rollout recovers success. A parallel cluster of
work makes rollout-based evaluation cheap and rigorous and thereby accepts its
cost~\citep{patil2026golden,snyder2026,vincent2024}; our study is the complement, testing
whether rollouts can be avoided at all.

What no prior work does is the full job on the offline side: measure the no-rollout gap across
regimes, compare a battery of offline metrics, test cross-task generalisation, frame it as
surrogate-endpoint calibration, and ship a reproducible protocol. The biomarker-qualification literature
\citep{prentice1989,buyse1998,joffe2009} gives the calibration vocabulary we borrow: a
surrogate must show both individual-level and trial-level association with the true endpoint
before anyone should trust it. That is precisely the question we put to validation loss.
